\documentclass[a4paper,12pt,twoside]{memoir}

% Castellano
\usepackage[spanish,es-tabla]{babel}
\selectlanguage{spanish}
\usepackage[utf8]{inputenc}
\usepackage[T1]{fontenc}
\usepackage{lmodern} % scalable font
\usepackage{microtype}
\usepackage{placeins}

\RequirePackage{booktabs}
\RequirePackage[table]{xcolor}
\RequirePackage{xtab}
\RequirePackage{multirow}

% Links
\usepackage[colorlinks]{hyperref}
\hypersetup{
	allcolors = {red}
}

% Ecuaciones
\usepackage{amsmath}

% Rutas de fichero / paquete
\newcommand{\ruta}[1]{{\sffamily #1}}

% Párrafos
\nonzeroparskip


% Imagenes
\usepackage{graphicx}
\newcommand{\imagen}[2]{
	\begin{figure}[!h]
		\centering
		\includegraphics[width=0.9\textwidth]{#1}
		\caption{#2}\label{fig:#1}
	\end{figure}
	\FloatBarrier
}

\newcommand{\imagenflotante}[2]{
	\begin{figure}%[!h]
		\centering
		\includegraphics[width=0.9\textwidth]{#1}
		\caption{#2}\label{fig:#1}
	\end{figure}
}



% El comando \figura nos permite insertar figuras comodamente, y utilizando
% siempre el mismo formato. Los parametros son:
% 1 -> Porcentaje del ancho de página que ocupará la figura (de 0 a 1)
% 2 --> Fichero de la imagen
% 3 --> Texto a pie de imagen
% 4 --> Etiqueta (label) para referencias
% 5 --> Opciones que queramos pasarle al \includegraphics
% 6 --> Opciones de posicionamiento a pasarle a \begin{figure}
\newcommand{\figuraConPosicion}[6]{%
  \setlength{\anchoFloat}{#1\textwidth}%
  \addtolength{\anchoFloat}{-4\fboxsep}%
  \setlength{\anchoFigura}{\anchoFloat}%
  \begin{figure}[#6]
    \begin{center}%
      \Ovalbox{%
        \begin{minipage}{\anchoFloat}%
          \begin{center}%
            \includegraphics[width=\anchoFigura,#5]{#2}%
            \caption{#3}%
            \label{#4}%
          \end{center}%
        \end{minipage}
      }%
    \end{center}%
  \end{figure}%
}

%
% Comando para incluir imágenes en formato apaisado (sin marco).
\newcommand{\figuraApaisadaSinMarco}[5]{%
  \begin{figure}%
    \begin{center}%
    \includegraphics[angle=90,height=#1\textheight,#5]{#2}%
    \caption{#3}%
    \label{#4}%
    \end{center}%
  \end{figure}%
}
% Para las tablas
\newcommand{\otoprule}{\midrule [\heavyrulewidth]}
%
% Nuevo comando para tablas pequeñas (menos de una página).
\newcommand{\tablaSmall}[5]{%
 \begin{table}
  \begin{center}
   \rowcolors {2}{gray!35}{}
   \begin{tabular}{#2}
    \toprule
    #4
    \otoprule
    #5
    \bottomrule
   \end{tabular}
   \caption{#1}
   \label{tabla:#3}
  \end{center}
 \end{table}
}

%
%Para el float H de tablaSmallSinColores
\usepackage{float}

%
% Nuevo comando para tablas pequeñas (menos de una página).
\newcommand{\tablaSmallSinColores}[5]{%
 \begin{table}[H]
  \begin{center}
   \begin{tabular}{#2}
    \toprule
    #4
    \otoprule
    #5
    \bottomrule
   \end{tabular}
   \caption{#1}
   \label{tabla:#3}
  \end{center}
 \end{table}
}

\newcommand{\tablaApaisadaSmall}[5]{%
\begin{landscape}
  \begin{table}
   \begin{center}
    \rowcolors {2}{gray!35}{}
    \begin{tabular}{#2}
     \toprule
     #4
     \otoprule
     #5
     \bottomrule
    \end{tabular}
    \caption{#1}
    \label{tabla:#3}
   \end{center}
  \end{table}
\end{landscape}
}

%
% Nuevo comando para tablas grandes con cabecera y filas alternas coloreadas en gris.
\newcommand{\tabla}[6]{%
  \begin{center}
    \tablefirsthead{
      \toprule
      #5
      \otoprule
    }
    \tablehead{
      \multicolumn{#3}{l}{\small\sl continúa desde la página anterior}\\
      \toprule
      #5
      \otoprule
    }
    \tabletail{
      \hline
      \multicolumn{#3}{r}{\small\sl continúa en la página siguiente}\\
    }
    \tablelasttail{
      \hline
    }
    \bottomcaption{#1}
    \rowcolors {2}{gray!35}{}
    \begin{xtabular}{#2}
      #6
      \bottomrule
    \end{xtabular}
    \label{tabla:#4}
  \end{center}
}

%
% Nuevo comando para tablas grandes con cabecera.
\newcommand{\tablaSinColores}[6]{%
  \begin{center}
    \tablefirsthead{
      \toprule
      #5
      \otoprule
    }
    \tablehead{
      \multicolumn{#3}{l}{\small\sl continúa desde la página anterior}\\
      \toprule
      #5
      \otoprule
    }
    \tabletail{
      \hline
      \multicolumn{#3}{r}{\small\sl continúa en la página siguiente}\\
    }
    \tablelasttail{
      \hline
    }
    \bottomcaption{#1}
    \begin{xtabular}{#2}
      #6
      \bottomrule
    \end{xtabular}
    \label{tabla:#4}
  \end{center}
}

%
% Nuevo comando para tablas grandes sin cabecera.
\newcommand{\tablaSinCabecera}[5]{%
  \begin{center}
    \tablefirsthead{
      \toprule
    }
    \tablehead{
      \multicolumn{#3}{l}{\small\sl continúa desde la página anterior}\\
      \hline
    }
    \tabletail{
      \hline
      \multicolumn{#3}{r}{\small\sl continúa en la página siguiente}\\
    }
    \tablelasttail{
      \hline
    }
    \bottomcaption{#1}
  \begin{xtabular}{#2}
    #5
   \bottomrule
  \end{xtabular}
  \label{tabla:#4}
  \end{center}
}



\definecolor{cgoLight}{HTML}{EEEEEE}
\definecolor{cgoExtralight}{HTML}{FFFFFF}

%
% Nuevo comando para tablas grandes sin cabecera.
\newcommand{\tablaSinCabeceraConBandas}[5]{%
  \begin{center}
    \tablefirsthead{
      \toprule
    }
    \tablehead{
      \multicolumn{#3}{l}{\small\sl continúa desde la página anterior}\\
      \hline
    }
    \tabletail{
      \hline
      \multicolumn{#3}{r}{\small\sl continúa en la página siguiente}\\
    }
    \tablelasttail{
      \hline
    }
    \bottomcaption{#1}
    \rowcolors[]{1}{cgoExtralight}{cgoLight}

  \begin{xtabular}{#2}
    #5
   \bottomrule
  \end{xtabular}
  \label{tabla:#4}
  \end{center}
}

%
% Nuevo comando para tablas internas
\newcommand{\innerTable}[1]{%
    \begin{tabular}[c]{m{0.7cm}|m{9.65cm}}
        Paso&  \\
        \hline
        #1
    \end{tabular} \\
}

\newcommand{\innerTables}[2]{% 
    \hline          Secuencia normal & \innerTable{#1}
    \hline          Excepciones & \innerTable{#2}
}

%
% Nuevo comando para tablas de requisitos funcionales
\newcommand{\outerTable}[9]{%
    \begin{table}[h!]
        \centering
        \caption{RF-#1 #2}
        \begin{tabular}[c]{m{2.2cm}|m{11cm}|}
            %\hline          RF-#1 & #2\\
            \hline\hline    Versión & #3 \\
            \hline          Autor & Rubén Marcos Gonzalez \\
            \hline          Tutor & Carlos Pardo Aguilar \\
            \hline          Objetivos asociados & #4 \\
            \hline          Requisitos asociados & #5 \\
            \hline          Descripción & #6 \\
            \hline          Precondición & #7 \\
            \innerTables#8
            \hline          Postcondicón & #9\\
            \hline
        \end{tabular}
        \label{RF:#1}
    \end{table}
}


\graphicspath{ {./img/} }

% Capítulos
\chapterstyle{bianchi}
\newcommand{\capitulo}[2]{
	\setcounter{chapter}{#1}
	\setcounter{section}{0}
	\chapter*{#2}
	\addcontentsline{toc}{chapter}{#2}
	\markboth{#2}{#2}
}

% Apéndices
\renewcommand{\appendixname}{Apéndice}
\renewcommand*\cftappendixname{\appendixname}

\newcommand{\apendice}[1]{
	%\renewcommand{\thechapter}{A}
	\chapter{#1}
}

\renewcommand*\cftappendixname{\appendixname\ }

% Formato de portada
\makeatletter
\usepackage{xcolor}
\newcommand{\tutor}[1]{\def\@tutor{#1}}
\newcommand{\course}[1]{\def\@course{#1}}
\definecolor{cpardoBox}{HTML}{E6E6FF}
\def\maketitle{
  \null
  \thispagestyle{empty}
  % Cabecera ----------------
\noindent\includegraphics[width=\textwidth]{cabecera}\vspace{1cm}%
  \vfill
  % Título proyecto y escudo informática ----------------
  \colorbox{cpardoBox}{%
    \begin{minipage}{.8\textwidth}
      \vspace{.5cm}\Large
      \begin{center}
      \textbf{TFG del Grado en Ingeniería Informática}\vspace{.6cm}\\
      \textbf{\LARGE\@title{}}
      \end{center}
      \vspace{.2cm}
    \end{minipage}

  }%
  \hfill\begin{minipage}{.20\textwidth}
    \includegraphics[width=\textwidth]{escudoInfor}
  \end{minipage}
  \vfill
  % Datos de alumno, curso y tutores ------------------
  \begin{center}%
  {%
    \noindent\LARGE
    Presentado por \@author{}\\ 
    en Universidad de Burgos --- \@date{}\\
    Tutor: \@tutor{}\\
  }%
  \end{center}%
  \null
  \cleardoublepage
  }
\makeatother


% Datos de portada
\title{título del TFG \\Documentación Técnica}
\author{Rubén Marcos González}
\tutor{Carlos Pardo Aguilar}
\date{\today}

\begin{document}

\maketitle



\cleardoublepage



%%%%%%%%%%%%%%%%%%%%%%%%%%%%%%%%%%%%%%%%%%%%%%%%%%%%%%%%%%%%%%%%%%%%%%%%%%%%%%%%%%%%%%%%



\frontmatter


\clearpage

% Indices
\tableofcontents

\clearpage

\listoffigures

\clearpage

\listoftables

\clearpage

\mainmatter

\appendix

\apendice{Plan de Proyecto Software}

\section{Introducción}

Se quiere crear un asistente de programación para C y en este anexo se comentará acerca del planificación que se ha seguido y su viabilidad

\section{Planificación temporal}

El proyecto se irá planificando semana a semana, estableciendo tareas y horario.
Si una tarea no se ha acabado en una semana se aplazará a la siguiente semana o se aplazará indefinidamente si han surgido muchas contrariedades en la realización de la misma dejando la opción de empezar con otras tareas.

\section{Estudio de viabilidad}

El proyecto en primera instancia pare plausible tanto técnicamente como legalmente pero no se sabe hasta que punto las diferencias que hay entre el lenguaje C y Python permitirán la realización del mismo.

\subsection{Viabilidad económica}

Considerando que ya hay asistentes de programación gratuitos en internet nuestra única opción económica es poner en manos de los usuarios donaciones para el proyecto como ya hace Eclipse.

\subsection{Viabilidad legal}

No hay ninguna oposición legal a este proyecto siempre que se declare la autoría tanto del parser\footnote{Enlace al repositorio del parser original:\url{https://github.com/eliben/pycparser}} como de los iconos\footnote{Fuente de los iconos:\url{https://www.flaticon.com/packs/essential-set-2}} utilizados en el proyecto.

\apendice{Especificación de Requisitos}

\section{Introducción}

El objetivo general de este programa es la depuración de código en lenguaje C. Para ello se requerirá poder abrir, crear, editar, compilar y depurar archivos con la extensión .c

\section{Objetivos generales}

\item OBJ-01:Depuración de código C.

\section{Catalogo de requisitos}

\item RI-01:Abrir archivos .c y .h
\item RI-02:Crear archivos .c y .h
\item RI-03:Editar archivos .c y .h
\item RI-04:Guardar archivos .c y .h
\item RI-05:Compilación de archivos .c
\item RI-06:Depuración de archivos .c

\section{Especificación de requisitos}



\outerTable{01}{Abrir archivos .c y .h}
{Versión 1.0}
{OBJ-01:Depuración de código C.}
{--}
{El programa deberá ser capaz de abrir archivos con la extensión .c y .h}
{El usuario deberá haber iniciado el asistente}
{{1: & solicita al sistema el explorador de archivos \\
\hline
2: & una vez seleccionado el archivo se accederá al editor \\}
{2: & el usuario abrirá un archivo con extensión no permitida \\
\hline
2: & el usuario abrirá un archivo corrupto \\}}
{El usuario deberá tener cargado un archivo con la extensión .c o .h}


\outerTable{02}{Crear archivos .c y .h}
{Versión 1.0}
{OBJ-01:Depuración de código C.}
{--}
{El programa deberá ser capaz de crear nuevos archivos con la extensión .c y .h}
{El usuario deberá haber iniciado el asistente}
{{1: & solicita al programador la ruta en la que quiere trabajar \\
\hline
2: & se solicita al programador el nombre del archivo \\
\hline
3: & se generará un archivo con el nombre y la extensión seleccionados \\}
{
1: & el usuario selecciona una ruta inexistente: se le preguntará sobre si crear las carpetas correspondientes o si quiere seleccionar una nueva ruta \\
\hline
3: & el usuario escogerá un nombre ya existente y se le preguntara si quiere sobreescribir el archivo \\}}
{El sistema deberá haber creado el archivo correspondiente y tenerlo cargado en el asistente}

\outerTable{03}{Editar archivos .c y .h}
{Versión 1.0}
{OBJ-01:Depuración de código C.}
{\item\nameref{RF:01}
\item\nameref{RF:02}}
{El programa deberá ser capaz de editar archivos con la extensión .c y .h}
{El usuario deberá tener cargado un archivo en el asistente}
{{1: & el usuario modificará lineas de código \\}
{
1: & el usuario no dispondrá de permisos de escritura por lo cual no se le permitirá modificar lineas de codigo \\}}
{El asistente deberá tener en memoria el archivo modificado para guardarlo posteriormente}

\outerTable{04}{Guardar archivos .c y .h}
{Versión 1.0}
{OBJ-01:Depuración de código C.}
{\nameref{RF:03}}
{El programa deberá ser capaz de guardar archivos con la extensión .c y .h}
{El usuario deberá tener cargado en el programa un archivo}
{{1: & solicita al sistema el explorador de archivos \\
\hline
2: & se sobrescribirá el archivo con el nombre y la extensión elegidas \\}
{2: & no haya espacio en el disco duro, se vuelve al paso 1 \\
\hline
2: & no existe dicho archivo, el sistema creará uno con el nombre y extensión elegidas \\}}
{El sistema deberá haber modificado correctamente el archivo}

\outerTable{05}{Compilar archivos .c}
{Versión 1.0}
{OBJ-01:Depuración de código C.}
{\nameref{RF:04}}
{El programa deberá ser capaz de compilar archivos .c y crear ejecutables a partir de ellos}
{El usuario deberá haber guardado los cambios del archivo}
{{1: & sobrescribirá el archivo .o correspondiente a dicho archivo \\
\hline
2: & sobrescribirá el ejecutable \\}
{1: & las modificaciones del archivo no han sido guardadas: se ejecutará el RF-04: Guardar archivos .c y .h \\
\hline
1: & el código de un error de compilación y se interrumpirá el proceso \\
\hline
1: & no existe el archivo .o correspondiente por lo que se creará uno nuevo \\
\hline
2: & no existe el ejecutable correspondiente por lo que se creará uno nuevo \\}}
{Disponer de un ejecutable funcional}

\outerTable{06}{Ejecutar linea de código (entrando en función)}
{Versión 1.0}
{OBJ-01:Depuración de código C.}
{\nameref{RF:05}}
{El programa deberá ser capaz de ejecutar una a una las lineas de código de un programa en lenguaje C}
{El usuario deberá haber compilado el archivo}
{{1: & se ejecutará una linea de código \\
\hline
2: & Irá mostrando las variables en memoria de ese programa acorde al \nameref{RF:09}\\}
{1: & no se dispone de un ejecutable funcional: se vuelve al \nameref{RF:05} \\
\hline
2: & encuentra una función por saltará al RF-07 Ejecutar linea de código (entrando en función) sobre esa función\\}}
{--}

\outerTable{07}{Ejecutar linea de código (saltando función)}
{Versión 1.0}
{OBJ-01:Depuración de código C.}
{\nameref{RF:05}}
{El programa deberá ser capaz de ejecutar una a una las lineas de código de un programa en lenguaje C}
{El usuario deberá haber compilado el archivo}
{{1: & se ejecutará una linea de código \\
\hline
2: & Irá mostrando las variables en memoria de ese programa acorde al \nameref{RF:09}\\}
{1: & no se dispone de un ejecutable funcional: se vuelve al \nameref{RF:05} \\}}
{--}

\outerTable{08}{Ejecutar código completo}
{Versión 1.0}
{OBJ-01:Depuración de código C.}
{\nameref{RF:05}}
{El programa deberá ser capaz de ejecutar un programa en lenguaje C completo}
{El usuario deberá haber compilado el archivo}
{{1: & se ejecutará el programa completo \\
\hline
2: & se mostrá el estado final de todas la variables acorde al \nameref{RF:09}\\}
{1: & no se dispone de un ejecutable funcional: se vuelve al \nameref{RF:05}\\}}
{--}

\outerTable{09}{Mostrar variables en memoria}
{Versión 1.0}
{OBJ-01:Depuración de código C.}
{\item\nameref{RF:06}
\item\nameref{RF:07}
\item\nameref{RF:08}}
{El programa deberá ser capaz de mostrar en una tabla las variables en memoria un un programa C}
{El usuario deberá empezado a ejecutar un programa C}
{{1: & cada vez que se cree o modifique una variable se mostrará en la tabla \\}
{-- & --\\}}
{--}
\apendice{Especificación de diseño}

\section{Introducción}

\section{Diseño de datos}

\section{Diseño procedimental}

\section{Diseño arquitectónico}



\apendice{Documentación técnica de programación}

\section{Introducción}

\section{Estructura de directorios}

\section{Manual del programador}

\section{Compilación, instalación y ejecución del proyecto}

\section{Pruebas del sistema}

\apendice{Documentación de usuario}

\section{Introducción}

\section{Requisitos de usuarios}

\section{Instalación}

\section{Manual del usuario}





\bibliographystyle{plain}
\bibliography{bibliografiaAnexos}

\end{document}
